% ============================================================
%  Section 7 -- Discussion
% ============================================================

\chapter{Discussion}

\section{Interprétation des résultats}

Cette étude montre que l'agent outillé domine le LLM unique sur les
dimensions fonctionnelles et opérationnelles. La différence vient de
l'utilisation de sources structurées : l'agent peut vérifier un statut serveur
ou retrouver une procédure au lieu de produire une réponse générique.

Ce résultat confirme l'intérêt des petits modèles lorsqu'ils sont couplés à des
outils. \code{phi3:latest} seul obtient une qualité faible, mais la même famille
de modèle placée dans un pipeline agentique passe le gate de release. Le gain ne
vient donc pas uniquement du modèle ; il vient de l'architecture.

\section{Limites}

Plusieurs limites doivent être explicitées.

\begin{itemize}
  \item \textbf{Dataset limité.} Les 21 tickets IT notés ne suffisent pas pour
  généraliser à tous les cas Michelin.
  \item \textbf{Configuration favorable à l'agent.} Le premier outil forcé et
  la synthèse déterministe réduisent la latence et les tokens. C'est utile pour
  tester le potentiel, mais cela doit être comparé à une configuration plus
  libre.
  \item \textbf{Juge LLM.} \code{gpt-5.2} apporte une évaluation scalable, mais
  il peut avoir ses propres biais, dont le biais de position et le biais de
  verbosité décrits dans les travaux LLM-as-a-judge \cite{llmjudge2023}. Une
  validation humaine reste nécessaire.
  \item \textbf{Artefacts hétérogènes.} Les campagnes A2A, Kaggle, MCP et
  Racing ne mesurent pas exactement le même objet.
  \item \textbf{Agents distants instables.} Quotas, timeouts et erreurs
  d'authentification affectent directement les résultats A2A.
\end{itemize}


\section{Perspectives}

Les perspectives prioritaires sont l'augmentation du dataset, l'annotation par
des experts support, la calibration du juge LLM, l'amélioration du RAG, la
génération automatique du rapport PDF depuis la CLI et l'ajout d'un tableau de
bord de suivi des régressions. Une autre piste consiste à router dynamiquement
les tickets : LLM unique pour les cas simples, agent outillé pour les cas
dépendants des sources, modèle plus fort pour les tickets critiques.
